\chapter{Conclusiones}

El Sistema AMR-FOD-Kitting para línea HTP A320 encaja con la asignatura de Sistemas Robóticos Móviles porque trabaja navegación interior, sensores, trazabilidad, seguridad, simulación y despliegue por flota. La propuesta no desplaza el proceso certificado de fabricación aeronáutica; actúa sobre logística interna, disponibilidad de útiles y prevención FOD.

La selección de un AMR omnidireccional responde a una restricción concreta: maniobrar en una nave industrial interior, cerca de estaciones, estructuras grandes y operarios. Frente a AGV, drones, orugas o robots con patas, la base omnidireccional reduce maniobras de aproximación y mantiene un coste compatible con un piloto.

La arquitectura separa percepción, navegación, planificación, control, supervisión, integración y base de datos. Con esa división puede empezar con 1 robot, pasar a 7 si los KPIs de seguridad y entrega se cumplen, y estudiar 17 robots cuando la flota demuestre disponibilidad, puntualidad y utilidad en rondas FOD. CoppeliaSim reduce riesgo antes de la implantación física y vincula la propuesta con los modelos estudiados en clase.

Comercialmente, el presupuesto queda desglosado por robot, sensores, software, ingeniería, instalación, formación, mantenimiento, logística, contingencia y margen. El ROI se calcula como recuperación de tiempo operativo, reducción de esperas, trazabilidad y apoyo al crecimiento de 40 a 80 HTP mensuales, sin prometer reducción directa de plantilla.
